\section{The gap hierarchy: what is protected, and by what}\label{sec:hierarchy}

The corpus states a ``hierarchy principle of structural persistence'': individual eigenvectors
are fragile, projectors are stabler, global invariants are stabler still, and a transport
should therefore be described by its projectors. Its later paper \Cc{} \cite{iqgt}, \S10.2
(``the tip and the tail), already states the law that makes this precise: a single
eigenvector's gap is the distance to its nearest neighbour, while the cluster projector's is
the distance from the cluster to the rest of the spectrum, so the two are governed by
different gaps and the subspace inherits the stability of the outer one. This section does
not re-derive that principle. It adds the constants: which rate, in which norm, with which
finite-pair denominator, and how much a rank-one projector gains --- nothing.

\subsection{Levels 0 and 2: eigenvalues and cluster projectors}

\begin{theorem}[Weyl]\label{thm:weyl}
Let $A,A'$ be self-adjoint with eigenvalues ordered increasingly. Then for every $i$,
$\abs{\lambda_i(A')-\lambda_i(A)}\le\norm{A'-A}$.
\Cc\ (Weyl). The constant $1$ is sharp: $A=\operatorname{diag}(0,\gamma)$,
$A'=A+\epsilon\,\sigma_1$ moves the eigenvalues by $\pm\epsilon$ to first order.
\end{theorem}

Eigenvalues therefore have no protection at all: any perturbation moves them, and at a crossing
they lose their identity. Projectors are different, and the difference is the gap.

\begin{theorem}[Projector stability: finite bound and its infinitesimal form]\label{thm:riesz}
Let $A,A'$ be self-adjoint on $\C^n$, let $\shape\subset\operatorname{spec}A$ be a spectral
cluster --- \emph{any} cluster, including one with spectrum on both sides of it --- and put
$\Gap:=\operatorname{dist}(\shape,\operatorname{spec}A\setminus\shape)>0$,
$\delta:=\norm{A-A'}$. Let $\PP,\PP'$ be the orthogonal projectors onto the cluster and onto the
part of $\operatorname{spec}A'$ enclosed by a fixed contour separating $\shape$ from the rest,
with $\delta<\Gap/2$ so that the count cannot change (Weyl). Then
\begin{enumerate}[label=(\alph*)]
\item \textbf{(finite pair)} $\norm{\PP-\PP'}\le\delta/(\Gap-\delta)$\label{eq:prb}. The
denominator is the \emph{union} gap: $\Gap-\delta\le\operatorname{dist}(\shape,
\operatorname{spec}A'\setminus\shape')$. In particular $\norm{\PP-\PP'}\le2\delta/\Gap$ on the
stable range, and the ratio $\norm{\PP-\PP'}/(\delta/\Gap)$ can exceed $1$: an adversarial search
reaches $1.078$ at $\delta/\Gap=0.346$ (\V, gate \gate{g04\_projector\_constant}).
\item \textbf{(infinitesimal)} along a $\C^1$ path $t\mapsto A(t)$ whose gap is at least $\Gap$,
\begin{equation}\label{eq:deriv}
   \norm{\dot\PP}\le\norm{\dot A}/\Gap ,
\end{equation}
and the constant $1$ is sharp: the family $A(t)=R(2t)\operatorname{diag}(1,-1)R(-2t)$ attains
it. This is the form used everywhere below; the finite-pair bound (a) is \emph{not} the
infinitesimal one with $\delta$ replaced by $\norm{\dot A}$.
\item \textbf{(angle)} for $\Theta:=\arcsin\norm{\PP-\PP'}\in[0,\pi/2]$,
$\Theta\le\arcsin\bigl(\delta/(\Gap-\delta)\bigr)$. Since $\delta/(\Gap-\delta)$ increases from
$0$ to $1$ across the stable range, the angle bound increases from $0$ to $\pi/2$ and
\emph{attains} $\pi/2$ at its endpoint $\delta=\Gap/2$, where the count of the cluster is about
to change. So $\pi/2$ is the exact constant of the angle bound on the closure of the stable
range, \emph{in the class of perturbations of arbitrary form} --- which is where it belongs, and
not in (a) or (b). Restricting the class sharpens it; see Remark~\ref{rem:tan}. The literature
states the estimate in this angle form: Albeverio and Motovilov \Cc{} prove
$\arcsin\norm{\PP-\QQ}\le M(\norm{V}/\Gap)$ for $\norm{V}<\Gap/2$, and their bound is strong
enough to \emph{ensure} the acute-angle case for $\norm{V}<c\,\Gap$ with
$c=0.454169\ldots<1/2$. That is a sufficient condition, and it says nothing about what
happens above $c\,\Gap$; in particular it does not assert that acuteness fails there. What
(c) supplies is the complementary information: the bound\'s own value tends to $\pi/2$ as
$\delta\uparrow\Gap/2$, so $\pi/2$ is the exact supremum of the angle bound over the stable
range. We do not claim that the bound is attained at $\delta=\Gap/2$ by an actual pair, and
an earlier draft of this paper incorrectly reported the threshold $c$ as evidence that it was.
\end{enumerate}
\Cc\ for (a): the Davis--Kahan $\sin\theta$ theorem in its union-gap form is imported, not
re-proved. \T\ for (b), which follows from (a) applied to the pair $(A(t),A(t+h))$ with
$\delta=\norm{\dot A}h$ and $\Gap-\delta\to\Gap$, divided by $h$; and for (c), which is
monotonicity of $\arcsin$. \V\ Sharpness of (b), and the failure of the constant-$1$ form of
(a): gate \gate{g04\_projector\_constant}.
\end{theorem}

\begin{remark}[A correction, and a correction of the correction]\label{rem:piover2}
This paper first quoted, as \Cc, a constant $\pi/2$ for the general (two-sided) case and $1$ for
the ordered case. Drafting the proof showed that the assertion ``$\norm{\PP-\PP'}\le
\delta/\Gap$ with constant $1$'' is \emph{false} in the normalization used here --- a targeted
search produces $\norm{\PP-\PP'}=0.3735$ against $\delta=0.3465$ --- because Davis--Kahan's
denominator is the union gap $\Gap-\delta$, not $\Gap$. Both errors are recorded rather than
replaced: the $\pi/2$ belongs to the angle bound, and the constant $1$ belongs to the
\emph{derivative}, which is what the hierarchy table uses. Whether the finite-pair bound improves for an ordered
cluster, and in which term, is not something we needed and do not claim.
\end{remark}

\begin{theorem}[Rank one, line-separated: the sharp off-diagonal bound]\label{thm:tan2}
Let $A$ be self-adjoint with a simple eigenvalue $a$ and unit eigenvector $u$, and suppose every
other eigenvalue of $A$ lies at least $d>0$ \emph{above} $a$. Let $V$ be self-adjoint and
\emph{off-diagonal} with respect to $H=\operatorname{span}u\oplus u^\perp$, put
$\delta:=\norm V$, and let $\nu$ be the eigenvector of $A+V$ for the eigenvalue near $a$. With
$\Theta$ the angle between the two lines,
\begin{equation}\label{eq:tan2}
   \abs{\tan 2\Theta}\le\frac{2\delta}{d},
   \qquad\text{hence}\qquad
   \tan\Theta\le\frac{\sqrt{d^{2}+4\delta^{2}}-d}{2\delta}<\frac{\delta}{d}\quad(\delta>0).
\end{equation}
Both inequalities are sharp: for $A=\operatorname{diag}(a+d,a)$ and $V=\delta\sigma_1$ one has
$\tan\Theta=\delta\big/\bigl(d/2+\sqrt{d^{2}/4+\delta^{2}}\bigr)$ and
$\abs{\tan 2\Theta}=2\delta/d$ \emph{exactly}, for every $\delta<d/2$.
\T\ Proof below; \V\ gate \gate{g06\_projector\_tan2} (12 assertions).
\end{theorem}

\begin{proof}
Write $\nu=\cos\Theta\,u+\sin\Theta\,w$ with $\norm w=1$, $w\perp u$, and put
$b:=\ip{Aw}{w}$, $\mu:=\ip{u}{Vw}=\ip{Vu}{w}$. Because $Au=au$, $Aw\perp u$, $Vu\perp u$ and
$Vw\perp w$ (the last two by off-diagonality), the eigenvalue equation $(A+V)\nu=\lambda\nu$ has
exactly two scalar consequences,
\[
   (\lambda-a)\cos\Theta=\mu\sin\Theta,\qquad
   (\lambda-b)\sin\Theta=\mu\cos\Theta .
\]
Dividing the first by $\cos\Theta$ and the second by $\sin\Theta$ and subtracting,
$b-a=\mu(\tan\Theta-\cot\Theta)=-2\mu/\tan 2\Theta$, so
$\abs{\tan 2\Theta}=2\abs\mu/\abs{b-a}$. Now $\abs\mu=\abs{\ip{Vu}{w}}\le\norm{Vu}\le\delta$,
and $b$, being a numerical-range point of the compression of $A$ to $u^\perp$, lies in the convex
hull of the rest of the spectrum, every point of which exceeds $a$ by at least $d$; hence
$b-a\ge d$. This gives \eqref{eq:tan2}'s first bound, and the second follows by solving
$2t/(1-t^{2})\le2\delta/d$ for $t=\tan\Theta<1$.
\end{proof}

\begin{remark}[Why line separation cannot be dropped, and what the convex-hull hypothesis adds]
\label{rem:numerical-range}
The proof uses $b-a\ge d$, and that can fail while $\operatorname{dist}(a,\operatorname{spec}A
\setminus\{a\})=d$ still holds: if the rest of the spectrum straddles $a$, the convex hull of the
rest contains $a$, the compression's numerical range can put $b$ at $a$, and the bound breaks ---
gate \gate{g06\_projector\_tan2}(d) exhibits a violation by $1.21$. So \emph{line separation} is
the hypothesis this argument needs, and it is weaker than what the literature's convex-hull
condition is for: that condition is what lets the argument be run for clusters of rank $k\ge2$,
where the scalar reduction below no longer closes because $\Theta$ is an operator angle and not
the angle of a single vector. The machinery for that case is the operator Riccati equation and
the direct rotation: existence and uniqueness of solutions of $A_+X-XA_-=XBX-B^*$ and the
$\tan\Theta$ theorem \Cc\ \cite{kmm1}, its $\tan 2\Theta$ generalisation \Cc\ \cite{kmm2}, and
the perturbation theory of spectra and spectral subspaces \Cc\ \cite{kmm3}. And note that \eqref{eq:tan2} is \emph{stronger} than the
Motovilov--Selin bound $\tan\Theta\le\delta/d$: for $\delta>0$ the right side of
\eqref{eq:tan2} is strictly smaller, so at rank one the sharp constant is the $\tan 2\Theta$ one.
\end{remark}

\begin{remark}[The angle constant depends on the class of perturbations]\label{rem:tan}
The bound (c) is sharp only over \emph{all} self-adjoint perturbations. If $V=A'-A$ is
\emph{off-diagonal} with respect to $\operatorname{Ran}\PP\oplus\operatorname{Ran}\QQ$, and the
convex hull of one spectral set misses the other, Motovilov and Selin \Cc\ prove the sharp
a priori $\tan\Theta$ theorem
\begin{equation}\label{eq:tan}
   \norm{\tan\Theta}\le\frac{\delta}{d},
   \qquad\text{equivalently}\qquad
   \norm{\PP-\PP'}\le\sin\Bigl(\arctan\frac{\delta}{\sqrt{d^{2}+\delta^{2}}}\Bigr),
\end{equation}
with $d$ the distance between the two spectral sets. Since $\arctan$ is bounded, the angle then
never reaches $\pi/2$ at any finite $\delta/d$, whereas \eqref{eq:prb}--(c) allows it at
$\delta=\Gap/2$. The two constants are therefore sharp in different classes: $\pi/2$ for the
unrestricted class, $\arctan$ for the off-diagonal one. At rank one with line separation
Theorem~\ref{thm:tan2} improves the second further, to $\tan 2\Theta\le 2\delta/d$. Both are recorded here rather than one
replacing the other, because the class is the hypothesis a reader has to check. \Cc\ for the
right-hand side of \eqref{eq:tan} --- the statement is Motovilov--Selin's Theorem~2 and we have not
re-derived it --- and \V\ for the transcription: gate \gate{g05\_tan\_theta} tests it over $300$
random off-diagonal perturbations of spectra with an interval cluster (worst ratio $0.987$ against
the bound), exhibits near-equality in the $2\times2$ model as $\delta/d\to0$ (ratio
$0.99999999$), and shows the hypothesis is load-bearing: a perturbation with a diagonal part
violates \eqref{eq:tan} by a factor $1.122$ at $\delta/d=0.318$. Motovilov--Selin's own sharpness
argument is not reproduced here.
\end{remark}

\begin{remark}[What the cluster's position does buy]\label{rem:ordered}
In Theorem~\ref{thm:riesz} the rate does not depend on where the cluster sits. What the position
does control is the \emph{angle scale}: in the examples of this section and of gate
\gate{g04}, a two-sided cluster at $\delta/\Gap\approx0.49$ rotates its subspace through an angle
about a third larger than an ordered one at the same $\delta$. An experiment that can arrange the
tracked cluster to be an end of the spectrum therefore buys a larger angular margin, not a better
constant.
\end{remark}

\subsection{Level 1: the correction --- a rank-one projector gains nothing}

\begin{theorem}[Rank one: no stability gain]\label{thm:rankone}
Let $v(t)$ be a $\C^1$ unit vector and $P=vv^*$. Then, with $\dot v_\perp:=\dot v-\ip{v}{\dot v}v$,
\begin{equation}\label{eq:rankone}
   \dot P=\dot v_\perp v^*+v\,\dot v_\perp^*,\qquad
   \norm{\dot P}_{\mathrm{op}}=\norm{\dot v_\perp},\qquad
   \norm{\dot P}_F=\sqrt2\,\norm{\dot v_\perp}.
\end{equation}
Consequently a rank-one projector is \emph{not} a more stable carrier than the eigenvector it
projects onto: in operator norm the two rates are equal, and in Frobenius norm the projector is
worse by $\sqrt2$.
\T\ (basis $(v,\dot v_\perp/\norm{\dot v_\perp})$, in which $\dot P$ is the block
$\bigl(\begin{smallmatrix}0&a\\a&0\end{smallmatrix}\bigr)$ with $a=\norm{\dot v_\perp}$).
\V\ Gate \gate{g02\_gap\_hierarchy} (a), 20 random pairs.
\end{theorem}

\begin{proof}
Write $\dot v=cv+\dot v_\perp$ with $c=\ip{v}{\dot v}\in i\R$ (since $\norm v=1$). Then
$\dot P=\dot vv^*+v\dot v^*=cvv^*+\dot v_\perp v^*+\bar cvv^*+v\dot v_\perp^*=\dot v_\perp v^*+v\dot v_\perp^*$
because $c+\bar c=0$. The two summands are orthogonal in the Frobenius inner product and each has
squared norm $\norm{\dot v_\perp}^2$, giving the Frobenius identity; on the two-dimensional space
spanned by $v$ and $\dot v_\perp$ the operator is $\bigl(\begin{smallmatrix}0&a\\a&0\end{smallmatrix}\bigr)$,
whose singular values are $a,a$.
\end{proof}

\begin{remark}[The correction, stated plainly]\label{rem:hierarchy-correction}
An earlier formulation of this paper, following the corpus, asserted
$\norm{\dot P}_F=\norm{\dot v_\perp}$. That is false by a factor $\sqrt2$; the identity holds in
operator norm. The correction matters because it removes the only reason to prefer a rank-one
projector: there is no gain, and in the norm most often computed there is a loss. \V\ Gate
\gate{g02\_gap\_hierarchy} (a) measures $\sqrt2$ to $10^{-12}$, and
\lean{FalseControls/} carries a deliberately false statement to the same effect.
\end{remark}

\subsection{The exception that proves the rule: a cluster through an exceptional point}

\begin{theorem}[Exceptional point and the trivial cluster projector]\label{thm:ep}
Let $B(t)=\bigl(\begin{smallmatrix}0&1\\t&0\end{smallmatrix}\bigr)$ for $t\in\C$. Then:
\begin{enumerate}[label=(\roman*)]
\item $\operatorname{spec}B(t)=\{\pm t^{1/2}\}$ for a branch of the square root; the eigenvalues
are continuous but not differentiable at $t=0$, and their splitting is $2\abs t^{1/2}$;
\item the unit eigenvectors are $(1,\pm t^{1/2})/\sqrt{1+\abs t}$; their distance is
$2\abs t^{1/2}/\sqrt{1+\abs t}$, so the eigenvectors \emph{coalesce} like $\abs t^{1/2}$ while
their derivatives diverge like $\abs t^{-1/2}$, and the condition number of the eigenvector
matrix is exactly $\abs t^{-1/2}$;
\item $B(t)$ is a single Jordan block only at $t=0$; for $t\neq0$ it is diagonalisable but its
eigenbasis is arbitrarily ill-conditioned as $t\to0$;
\item the Riesz projector onto the pair $\{\pm t^{1/2}\}$ is the identity, hence analytic, and
its derivative vanishes.
\end{enumerate}
\T\ (direct computation). \V\ Gate \gate{g02\_gap\_hierarchy} (b), residuals below $10^{-9}$ for
$\abs t$ down to $10^{-6}$.
\end{theorem}

\begin{remark}[What the exceptional point does and does not destroy]
\label{rem:ep-scope}
The singular behaviour at $t=0$ is in the \emph{basis}, not in the \emph{cluster}: the pair's
projector is analytic, the pair's trace is analytic, and every quantity that depends only on the
cluster is analytic through the degeneracy. Individual spectral projectors, eigenvectors and
condition numbers diverge as $\abs t^{-1/2}$; that is the price of insisting on individual
levels. So the corpus's phrase ``exceptional-point regularisation by Riesz projectors'' is
correct, but only in this precise sense: the regularisation is available for a cluster whose
\emph{external} gap stays open, and it regularises nothing about the individual levels inside
the cluster. Where the coupling between the exceptional block and the rest of the spectrum is
nonzero, the cluster projector is no longer constant, and the constant governing its derivative
is the external gap $\Gap_{\mathrm{ext}}$; that coupled statement is \V\ for the families in
\gate{results/swarm/D5\_m4.md} but is not proved here, and it is listed as an open problem
(\S\ref{sec:open}, O3).
\end{remark}

\subsection{Level 3: the integers, and what they count}

\begin{theorem}[Winding, and why it is not a crossing count]\label{thm:winding}
Let $t\mapsto H(t)=a(t)\sigma_1+b(t)\sigma_2+c(t)\sigma_3$ be continuous on the circle
$t\in S^1$.
\begin{enumerate}[label=(\alph*)]
\item If $a+ib$ never vanishes, its winding number $w\in\mathbb{Z}$ is defined and is invariant
under every continuous deformation keeping $a+ib\neq0$.
\item The winding is \emph{not} the spectral flow. $H(t)=\cos t\,\sigma_1+\sin t\,\sigma_2$ has
$w=1$ and no eigenvalue crossing any level, so its spectral flow is $0$; and $H=\sigma_3$ has
$a+ib\equiv0$, so its winding is undefined although its spectrum is gapped at every level
$0<\mu<1$.
\item If $c\equiv0$, the lower band's Berry phase $\gamma_-$ satisfies
$\exp(i\gamma_-)=(-1)^w$.
\end{enumerate}
\T\ (a) is the argument principle; (c) is the standard two-level computation
$\gamma_-=\pi$ for $w=1$, verified independently for a real and for a complex Hermitian loop.
\V\ Gate \gate{g02\_gap\_hierarchy} (c) and \gate{results/swarm/check\_D6\_m8.py}.
\end{theorem}

\begin{remark}[The corrected currency]\label{rem:currency}
An earlier formulation said ``spectral flow equals winding''. Theorem~\ref{thm:winding} shows
that neither direction survives as stated. What survives is the useful half: when a two-band
family never degenerates on the loop, the winding is an integer, it is a homotopy invariant of
the gapped family, and it obstructs a globally continuous gapped trivialisation. That is what a
``topological currency'' is, and it is a homotopy class rather than a count of crossings.
\end{remark}

\subsection{The hierarchy, stated as a theorem's worth of content}

\begin{center}
\begin{tabular}{@{}p{0.6cm}p{3.4cm}p{3.4cm}p{4.2cm}@{}}
\toprule
level & invariant & protected by & failure mode \\
\midrule
0 & eigenvalue & nothing & Weyl: moves by $\norm{\dot A}$; crossing destroys the label \\
1 & eigenvector, or a rank-one projector & the internal gap $\Gap_{\mathrm{int}}$ &
\emph{no gain} (\S\ref{sec:hierarchy}.2): rate $\norm{\dot A}/\Gap_{\mathrm{int}}$, $\sqrt2$ worse in Frobenius; at an EP diverges like $\abs t^{-1/2}$ \\
2 & cluster projector & the external gap $\Gap_{\mathrm{ext}}$ &
analytic through interior degeneracies; rate $\norm{\dot A}/\Gap_{\mathrm{ext}}$, constant $1$ in the infinitesimal form \\
3 & trace, rank, winding, spectral flow & the external gap, and integrality & locally constant; jumps only in integers \\
\bottomrule
\end{tabular}
\end{center}

\begin{remark}[The principle, corrected]\label{rem:principle}
\emph{A transport invariant is stable exactly to the extent that the gap protecting it is open,
and the relevant gap is the one across the boundary of the cluster that defines the invariant ---
never the gap inside it.} The principle is the corpus's own (\cite{iqgt}, \S10.2). What is added
here is the quantitative row: passing from an eigenvector to its \emph{rank-one} projector buys
nothing --- the operator-norm rate is identical and the Frobenius rate is worse by $\sqrt2$ ---
because they are the same object; passing from a level to the \emph{cluster} containing it buys
the external gap, whose protection Theorem~\ref{thm:riesz} makes explicit. An experiment that wants stability should
therefore enlarge its cluster until the gap it controls is the outer one, and a theory that wants
to claim stability should name that gap.
\end{remark}
